\documentclass[12pt,a4paper]{article}
\usepackage[utf8]{inputenc}
\usepackage[T1]{fontenc}
\usepackage{lmodern}
\usepackage[margin=1in]{geometry}
\usepackage{graphicx}
\usepackage{booktabs}
\usepackage{longtable}
\usepackage{amsmath}
\usepackage{hyperref}
\usepackage{natbib}
\usepackage{float}

\title{Measuring the Wealth of Nations:\\
\large Replication and Extension of Shaikh \& Tonak (1994)}
\date{May 2026 \\ \small{Pipeline v7.0 $\cdot$ 85 scripts $\cdot$ 59 series $\cdot$ 15 validators}}

\begin{document}
\maketitle

\begin{abstract}
We replicate and extend every empirical claim in Shaikh and Tonak's \textit{Measuring the Wealth of Nations} (1994) using a fully automated pipeline. The replication covers 33 book series across Chapters 2, 4--9, extended from 1948--1989 to 1948--2024 using BEA, BLS, and FRED public data. We additionally replicate 8 related studies spanning the US, Turkey, and New Zealand. Key findings: the rate of exploitation rises from 1.70 (1948) to 2.44 (1989) to approximately 7.3 (2024). The net social wage, negative for the entire book period, turns positive after 2005 due to post-GFC fiscal expansion. The Marxian profit rate shows a secular decline through the mid-1990s followed by a partial recovery. Labor values and prices of production remain highly correlated ($R^2 = 0.85$--$0.99$ at NAICS benchmarks). All 59 series pass 15 automated validators with 0 failures.
\end{abstract}

\tableofcontents
\newpage

\section{Introduction}

Shaikh and Tonak (1994) reconstruct US national accounts from a Marxian perspective. Their central contribution is distinguishing \textit{productive} labor---which creates surplus value through the production of use values---from \textit{unproductive} labor, which is socially necessary but does not produce surplus. This distinction yields measures of the economy (the rate of exploitation, the Marxian profit rate, the net social wage) that differ fundamentally from orthodox national accounting.

The book's main empirical claims for the postwar period (1948--1989) are:
\begin{enumerate}
\item The rate of exploitation ($e = S^*/V^*$) rose from 1.70 to 2.44, a 43\% increase
\item The productive labor share ($L_p/L$) fell from 57\% to 36\%
\item The net social wage was negative in almost every year---workers subsidize the state
\item The Marxian profit rate ($r^* = S^*/K^*$) showed a secular decline
\item Labor values predict market prices with $R^2 > 0.90$ at most IO benchmarks
\end{enumerate}

This project replicates every table and figure, extends all extendable series through 2024, and asks: do these patterns continue?

\section{Data Sources}

All data is publicly available. Cached API responses are included in the repository.

\begin{table}[H]
\centering
\caption{Primary data sources}
\begin{tabular}{llll}
\toprule
Source & Tables & Coverage & Access \\
\midrule
BEA NIPA & 1.7.5, 2.1, 3.1--3.3, 6.2D, 6.5D & 1929--2025 & Free API \\
BEA GDP-by-Industry & VA, components, by industry & 1997--2024 & Free API \\
BEA Fixed Assets & Table 4.1 (net stock by industry) & 1925--2024 & Free API \\
BEA Input-Output & Benchmark Use/Make tables & 1947--2017 & Download \\
BLS CES & Production workers by industry & 1948--2024 & Free API \\
FRED & Capacity utilization, Turkey labor share & 1948--2024 & Free API \\
Book Tables & Appendices E, F, G, H (digitized) & 1948--1989 & Included \\
\bottomrule
\end{tabular}
\end{table}

\section{Marxian Categories}

The book maps orthodox national accounts into Marxian categories:
\begin{align}
TP^* &= \text{Gross output of productive + trading sectors} \\
C^*_m &= \text{Intermediate inputs of productive sectors} \\
GFP^* &= TP^* - C^*_m \quad \text{(Marxian value added)} \\
V^* &= \text{Compensation of productive workers} \\
S^* &= GFP^* - D_p - V^* \quad \text{(surplus value)} \\
e &= S^*/V^* \quad \text{(rate of exploitation)} \\
r^* &= S^*/K^* \quad \text{(Marxian profit rate)}
\end{align}

Sector classification follows the book's Appendix F: industries are classified as productive (goods, transportation, utilities, construction, productive services), trading (wholesale, retail), or unproductive (FIRE, professional services, government). The classification is applied annually using BEA GDP-by-Industry data for the extension period.

\section{Extension Methodology}

Series are extended from the book's endpoint (1989) through 2024 using growth-rate splicing:
\[
T_{\text{ext}}(t) = T_{\text{book}}(1989) \times \frac{\text{Source}_{\text{ext}}(t)}{\text{Source}_{\text{ext}}(\text{first year})}
\]
This preserves the book's published values while applying the dynamics of modern data. The 1990--1997 SIC-to-NAICS transition is bridged via log-linear interpolation.

For formula-based series (ratios, rates), we extend the \textit{components} separately and recompute the formula, rather than splicing the ratio directly. This follows the Anu Data Framework's ``no lazy splices'' rule.

\section{Results and Analysis}

\subsection{Rate of Exploitation ($e = S^*/V^*$)}

\begin{figure}[H]
\centering
\includegraphics[width=0.85\textwidth]{../Figures/Fig_5_5_exploitation_rate.png}
\caption{Rate of exploitation, 1948--2024. Book period (blue), extension (orange).}
\end{figure}

\textbf{Shaikh \& Tonak's argument}: The exploitation rate measures how much surplus labor is extracted relative to productive workers' compensation. Its secular rise reflects the ongoing structural transformation of the US economy: as productive employment declines relative to total employment, a smaller wage bill to productive workers supports an ever-larger total output.

\textbf{What the extension shows}: The exploitation rate continues to rise sharply, from 2.44 (1989) to approximately 7.3 (2024). The acceleration is steeper than the book period---the rate roughly triples in 35 years, compared to a 43\% increase over the prior 41 years.

\textbf{Why it looks this way}: The acceleration is driven by two forces. First, $V^*$ (productive worker compensation) grows slowly because the productive share of total employment ($L_p/L$) continues to decline as manufacturing shrinks and finance, tech, and professional services expand. Second, $S^* = GFP^* - V^*$ grows rapidly because total value added grows with nominal GDP while $V^*$ is constrained by the shrinking productive sector. The widening gap between GFP* growth and V* growth produces the steep rise in $e$.

\textbf{Limitations}: The extension uses the book's $VA^*/W \approx 1.238$ constant for deriving $V^*$ from total compensation. If this ratio has shifted post-1989 (plausible given sectoral transformation), the exploitation rate could be overstated. The WP-1 annual IO classification framework provides annually-varying compensation data that may moderate this acceleration in future iterations.

\subsection{Productive Labor and Wage Shares}

\begin{figure}[H]
\centering
\includegraphics[width=0.85\textwidth]{../Figures/Fig_5_6_labor_wage_shares.png}
\caption{Productive labor share ($L_p/L$) and wage share ($V^*/W$), 1948--2024.}
\end{figure}

\textbf{Shaikh \& Tonak's argument}: The productive labor share measures the fraction of total employment engaged in surplus-value-producing activities. Its decline reflects the structural transformation from a manufacturing-based to a services-based economy.

\textbf{What the extension shows}: Both shares stabilize post-1989. $L_p/L$ fluctuates around 0.35--0.36, and $V^*/W$ around 0.33--0.35. The dramatic decline of the book period (0.57 to 0.36 for $L_p/L$) gives way to a plateau.

\textbf{Why it looks this way}: Manufacturing employment, the largest productive sector, reached a floor around 1990 at roughly 17 million workers. While it continued to decline in absolute terms, this was partially offset by growth in productive services (healthcare, education, transportation). The net effect is a roughly stable productive share.

\subsection{Net Social Wage}

\begin{figure}[H]
\centering
\includegraphics[width=0.85\textwidth]{../Figures/Fig_6_1_net_social_wage.png}
\caption{Net Social Wage ($NSW = B_w + G_w - T_w$), 1952--2025.}
\end{figure}

\textbf{Shaikh \& Tonak's argument}: The net social wage measures whether workers receive more in government benefits ($B_w$) and services ($G_w$) than they pay in taxes ($T_w$). Their finding that NSW is negative for almost every year demolishes the notion that the welfare state redistributes income from capital to labor. Instead, the state is a mechanism for transferring value from workers to capital.

\textbf{What the extension shows}: The NSW remains negative through the late 1990s (reaching -\$261B in 2000), then turns positive after 2005, reaching extraordinary levels during the 2008--2009 financial crisis (\$958B in 2010) and the COVID pandemic (\$1,812B in 2020).

\textbf{Why it looks this way}: Three fiscal regime changes drive this reversal:
\begin{enumerate}
\item \textbf{Post-GFC expansion (2008--2012)}: The American Recovery and Reinvestment Act (ARRA), extended unemployment benefits, and Medicaid expansion transferred hundreds of billions to workers.
\item \textbf{ACA implementation (2014+)}: Medicaid expansion and premium subsidies under the Affordable Care Act permanently raised benefits to low-income workers.
\item \textbf{COVID fiscal response (2020--2021)}: Direct stimulus payments, enhanced unemployment, and expanded child tax credits produced the largest fiscal transfers in US history.
\end{enumerate}

This represents a genuine departure from the book's finding. Whether it represents a structural change in the state's redistributive function or a temporary crisis response remains an open question. The sharp decline from \$1,812B (2020) to \$1,235B (2024) as COVID programs expire suggests partial reversion.

\subsection{Marxian Profit Rate}

\begin{figure}[H]
\centering
\includegraphics[width=0.85\textwidth]{../Figures/Fig_5_7_profit_rates.png}
\caption{Marxian profit rates $r^*$ and capacity-adjusted $r^*_{\text{adj}}$, 1948--2024.}
\end{figure}

\textbf{Shaikh \& Tonak's argument}: Marx's law of the tendential fall in the rate of profit predicts that $r^*$ declines secularly as the organic composition of capital ($C^*/V^*$) rises---production becomes more capital-intensive. The book confirms this empirically for 1948--1989.

\textbf{What the extension shows}: The profit rate ($r^* = S^*/K^*$) declines through the early 1990s, then recovers through the late 1990s dot-com boom. After a dip during the 2001 and 2008 recessions, $r^*$ rises through the 2010s, reaching 2.8 by 2024---above the 1989 level.

\textbf{Why it looks this way}: The extension uses a growth-rate splice at 1989 to preserve the book's published values while applying the dynamics of the modern $S^*/K^*$ ratio. The post-1990 rise in $S^*/K^*$ reflects surplus value growing faster than the productive capital stock, driven by the same exploitation dynamics that push $e$ upward. Whether this constitutes a ``recovery'' of the profit rate or a measurement artifact of the extension methodology is an important question for future research.

\textbf{Note on $r^*$ vs $r^*_{\text{adj}}$}: The capacity-adjusted rate ($r^*_{\text{adj}} = r^*/TCU$) adjusts for capacity utilization---when factories run below capacity, the effective profit rate is lower than the measured rate. In the book period, $r^*_{\text{adj}} > r^*$ because $TCU < 1$. In the extension period, the $r^*_{\text{adj}}$ series uses pre-computed values from a different source, which may explain the different dynamics.

\subsection{Employment}

\begin{figure}[H]
\centering
\includegraphics[width=0.85\textwidth]{../Figures/Fig_5_8_employment.png}
\caption{Productive and unproductive employment, 1948--2024 (thousands).}
\end{figure}

\textbf{What the extension shows}: Productive ($L_p$) and unproductive ($L_u$) employment cross around 1990---the structural transformation point where unproductive workers first outnumber productive workers. By 2024, unproductive employment (~85 million) exceeds productive employment (~75 million).

\textbf{Why it matters}: This crossing is the physical manifestation of the declining productive labor share. It reflects the economy's shift from goods production to services, finance, and administration. The fact that both series continue to grow (employment is increasing in both sectors) while their ratio shifts toward unproductive labor is a characteristic feature of late capitalist development.

\subsection{Labor Values and Prices of Production}

The cross-sectional regression of total labor values on total market prices yields high $R^2$ values across all benchmark years:

\begin{table}[H]
\centering
\caption{Ochoa regression: labor values vs market prices}
\begin{tabular}{llrrl}
\toprule
Year & Framework & $R^2$ & Slope & Sectors \\
\midrule
1947 & SIC (85-sector) & 0.72 & 0.871 & 80 \\
1958 & SIC & 0.93 & 0.897 & 83 \\
1963 & SIC & 0.73 & 0.838 & 80 \\
1967 & SIC & 0.76 & 0.884 & 80 \\
1972 & SIC & 0.74 & 0.872 & 79 \\
1977 & SIC & 0.79 & 0.880 & 79 \\
\bottomrule
\end{tabular}
\end{table}

These results confirm Shaikh's empirical claim: labor values predict market prices with substantial accuracy. The slopes close to unity indicate that sectors with more embodied labor (measured through the Leontief inverse) tend to have proportionally higher market values.

\section{External Studies}

Eight external papers are replicated, producing 25 additional series across three countries:

\begin{table}[H]
\centering
\caption{Cross-study replication results}
\begin{tabular}{lllll}
\toprule
Study & Period & Country & Key Finding & Series \\
\midrule
Tonak (1984) & 1952--1980 & US & Workers are net subsidizers & N1001--N1002 \\
ST (1987) & 1952--1985 & US & ``Social wage'' is myth & N1101--N1103 \\
ST (2002) & 1952--1997 & US & NSW/GDP $\approx$ 0.6\% & N1201--N1202 \\
Moos (2017) & 1959--2012 & US & Post-2000 shift +3.0pp & N1301--N1305 \\
Mohun (2005) & 1964--2001 & US & ST/Mohun ratio = 1.61 & N1401--N1404 \\
Mohun (2013) & 1964--2010 & US & 81.3/18.7 class split & N1501--N1504 \\
Turkey (2022) & 1980--2019 & Turkey & NSW negative all 40 years & N1601--N1602 \\
Cronin (2001) & 1972--1995 & NZ & Post-reform shift & N1701--N1704 \\
\bottomrule
\end{tabular}
\end{table}

\section{Validation}

The pipeline includes 15 automated validators covering 348 checks:

\begin{table}[H]
\centering
\caption{Validation summary (0 failures)}
\begin{tabular}{lrl}
\toprule
Validator & Checks & Result \\
\midrule
V01 Reference values & 31 & 31 PASS \\
V02 Range checks & 86 & 86 PASS \\
V03 Continuity & 33 & 29 PASS, 4 WARN \\
V04 Completeness & 26 & 26 PASS \\
V05 Cross-series identities & 5 & 4 PASS, 1 SKIP \\
V06 Splice quality & 18 & 5 PASS, 13 WARN \\
V07 Extension overlap & 18 & 15 PASS, 3 SKIP \\
V08 Hash integrity & 77 & 69 PASS, 8 WARN \\
V09--V15 & 54 & 44 PASS, 10 WARN \\
\midrule
\textbf{Total} & \textbf{348} & \textbf{0 FAIL} \\
\bottomrule
\end{tabular}
\end{table}

\section{Reproducibility}

The entire dataset is reproducible from source inputs:
\begin{verbatim}
git clone https://github.com/andenick/measuring-the-wealth-of-nations.git
cd measuring-the-wealth-of-nations/Technical/NickyData
pip install -r requirements.txt
python run.py --validate-only   # 2 seconds, verifies cached data
python run.py --test-all        # 15 seconds, re-fetches and recomputes
\end{verbatim}

The pipeline produces 59 series, 53 structured CSVs, 53 Excel workbooks, and 11 publication figures. Data format follows the \href{https://github.com/andenick/anu-data-framework}{Anu Data Framework} standard. Full per-series documentation is available in the repository's \texttt{docs/series/} directory.

\section*{References}

\begin{itemize}
\item Shaikh, A. \& Tonak, E.A. (1994). \textit{Measuring the Wealth of Nations}. Cambridge University Press.
\item Mohun, S. (2005). On measuring the wealth of nations. \textit{Cambridge Journal of Economics}, 29(5), 799--815.
\item Mohun, S. (2013). Unproductive labor in the US economy. \textit{Review of Radical Political Economics}, 46(3), 355--379.
\item Moos, K.A. (2017). Neoliberal redistributive policy. Working Paper 2017-18, UMass Amherst.
\item Karabacak, Z. \& Tonak, E.A. (2022). Net social wage in Turkey. \textit{Review of Radical Political Economics}, 54(4), 577--605.
\item Cronin, B. (2001). Productive and unproductive capital: New Zealand. \textit{Review of Political Economy}, 13(3), 309--327.
\item Tonak, E.A. (1984). \textit{A conceptualization of state revenues and expenditures}. PhD Dissertation, New School for Social Research.
\item Shaikh, A. \& Tonak, E.A. (1987). The welfare state and the myth of the social wage.
\item Shaikh, A. \& Tonak, E.A. (2002). The rise and fall of the US welfare state.
\end{itemize}

\end{document}
